\section{Conclusion}
\label{sec:conclusion}

Legal artificial intelligence has evolved from early rule-based systems and case-based reasoning toward
domain-adapted language models, retrieval-augmented generation, and increasingly sophisticated legal
information workflows. Each stage has addressed limitations of earlier approaches, but the field's
development also reveals a persistent distinction between producing plausible legal language and providing
a legally supportable answer.

Early systems demonstrated that legal reasoning could be represented through explicit rules, structured
knowledge, and comparisons between cases. These approaches offered interpretability but faced challenges
in maintaining extensive legal knowledge and adapting to changing legal materials. Transformer-based
language models subsequently improved the processing of legal language, while benchmarks such as
LEGAL-BERT-related evaluations, LexGLUE, and LegalBench enabled more systematic assessment of specialized
language understanding and legal reasoning. Nevertheless, benchmark performance alone cannot establish
that a system retrieves current authority, interprets it correctly, or cites it accurately.

Retrieval-augmented generation introduced an important architectural shift by connecting language
generation to external evidence. Legal RAG systems can retrieve statutory provisions, judgments, and other
source materials at inference time, making answers more traceable than approaches that depend exclusively
on model parameters. Research on precise legal passage retrieval, structured document drafting, and
jurisdiction-specific question answering illustrates the breadth of this development. Yet RAG does not
guarantee reliability: relevant passages may be omitted, retrieved fragments may lose essential context,
and generated claims may exceed what the evidence supports.

The comparative analysis identifies three findings. First, retrieval quality, legal correctness, and
citation support are distinct dimensions and must be evaluated separately. Second, legal document
structure, jurisdiction, and temporal validity are essential to determining whether retrieved evidence is
applicable. Third, more complex architectures, including knowledge graphs and agentic workflows, offer
opportunities to address these problems, but their value must be demonstrated through controlled
evaluation rather than inferred from architectural sophistication.

These findings are particularly relevant to multilingual and jurisdiction-specific legal information
retrieval. In contexts such as Pakistan, effective systems must connect user queries with the appropriate
legal sources while preserving statutory hierarchy, judicial context, and source provenance. Supporting
English, Urdu, and Roman Urdu adds further challenges involving terminology alignment, transliteration,
and cross-lingual semantic matching. A system that handles these challenges successfully could improve
access to legal information, but only if its outputs remain traceable to applicable authorities and its
uncertainty is handled responsibly.

Future research should therefore prioritize structure-preserving retrieval, provenance-aware legal
knowledge graphs, multilingual query processing, claim-level citation verification, and agentic workflows
with explicit quality gates. Equally important are reproducible benchmarks that test the complete
evidence-to-answer pipeline, including difficult cases involving exceptions, conflicting authorities,
temporal changes, and insufficient evidence. Such evaluation should report retrieval effectiveness, legal
correctness, citation support, latency, and computational cost without collapsing them into a single
score.

Ultimately, the goal of legal AI should not be to replace legal judgment with fluent automated responses.
It should be to make legal information more accessible, systematic, and auditable while preserving the
distinction between retrieved evidence, model-generated interpretation, and conclusions requiring
professional review. The next stage of legal AI will depend less on generation alone and more on the
ability to retrieve, contextualize, verify, and reproduce evidence-supported legal reasoning.
